% 妊神星（综述）
% license CCBYSA3
% type Wiki

本文根据 CC-BY-SA 协议转载翻译自维基百科\href{https://en.wikipedia.org/wiki/Haumea}{相关文章}。

帕洛玛天文台是一座位于美国加利福尼亚州圣迭戈县帕洛玛山的天文研究台站。该天文台由加州理工学院（Caltech）拥有并运营。观测时间由加州理工学院及其科研合作单位共享，这些单位包括喷气推进实验室（JPL）、耶鲁大学$^{[1]}$、以及中国科学院国家天文台$^{[2]}$。

该天文台运行多台望远镜，包括 200 英寸（5.1 米）的黑尔望远镜$^{[3]}$，48 英寸（1.2 米）的塞缪尔·奥申望远镜$^{[4]}$（用于 Zwicky 暂现天体设施 ZTF）$^{[5]}$，帕洛玛 60 英寸（1.5 米）望远镜$^{[6]}$，以及 30 厘米（12 英寸）的 Gattini-IR 望远镜$^{[7]}$。已退役的仪器包括帕洛玛试验台干涉仪，以及天文台最早启用的望远镜——建于 1936 年的 18 英寸（46 厘米）施密特相机。
\subsection{历史}
\begin{figure}[ht]
\centering
\includegraphics[width=6cm]{./figures/15c1f9cc6c3cbd5f.png}
\caption{帕洛玛山天文台出现在 1948 年的美国邮票上} \label{fig_Haumea_1}
\end{figure}
\subsubsection{Hale 对大型望远镜与帕洛玛天文台的愿景}
天文学家乔治·埃勒里·黑尔（George Ellery Hale），其愿景促成了帕洛玛天文台的建立，曾连续四次建造出当时世界上最大口径的望远镜$^{[8]}$。他在 1928 年发表了一篇文章，提出了后来成为 200 英寸帕洛玛反射望远镜的构想；这是一份向美国公众发出的邀请，旨在让他们了解大型望远镜如何帮助解答宇宙本质的基本问题。黑尔随后于 1928 年 4 月 16 日致信洛克菲勒基金会的国际教育委员会（后来并入一般教育委员会），在信中提出对此项目的资金申请。在信中，黑尔写道：

“在推动科学进步方面，没有任何方法能像开发全新的、更强大的仪器和研究方法那样富有成效。一台更大的望远镜不仅能够在空间穿透力与摄影分辨率上提供必要的提升，还可以让源自物理与化学最新基础性进展的理念与装置得以应用。”
\subsubsection{黑尔望远镜}
这台 200 英寸望远镜以天文学家兼望远镜制造者乔治·埃勒里·黑尔的名字命名。它由加州理工学院利用洛克菲勒基金会提供的 600 万美元资助建造，采用的是康宁玻璃厂在乔治·麦考利（George McCauley）指导下制造的 Pyrex 镜坯。J.A. Anderson 博士是最初的项目经理，其任命可追溯至 1930 年代早期$^{[9]}$。该望远镜（当时世界上最大）于 1949 年 1 月 26 日首次观测（first light），目标为 NGC 2261$^{[10]}$。美国天文学家埃德温·鲍威尔·哈勃（Edwin Powell Hubble）是第一位使用该望远镜的天文学家。

这台 200 英寸望远镜自 1949 年起一直保持世界最大望远镜的地位，直到 1975 年俄罗斯 BTA-6 望远镜首次观测。使用黑尔望远镜的天文学家在宇宙尺度上发现了类星体（后被归入所谓的活动星系核的一类）。他们研究了恒星族群的化学性质，从而理解了恒星核合成在宇宙中按观测丰度产生元素的起源，并已发现数千颗小行星。在纽约康宁市的康宁社区学院，存放着一台 1/10 比例的望远镜工程模型（康宁玻璃厂 — 即现今的康宁公司所在地），该模型曾用于发现至少一颗小行星——34419 Corning。
\subsubsection{建筑与设计}
\begin{figure}[ht]
\centering
\includegraphics[width=6cm]{./figures/6fd368c863632c16.png}
\caption{黑尔望远镜穹顶} \label{fig_Haumea_2}
\end{figure}
拉塞尔·W·波特（Russell W. Porter）设计了天文台建筑的装饰艺术（Art Deco）风格，包括 200 英寸黑尔望远镜的穹顶。波特还负责了黑尔望远镜与施密特相机的大部分技术设计工作，绘制了一系列横截面工程图。波特与许多工程师及加州理工学院委员会成员合作进行这些设计$^{[11]}$ $^{[12]}$ $^{[13]}$。
马克斯·梅森（Max Mason）负责施工，而西奥多·冯·卡门（Theodore von Karman）参与了工程设计。
\subsubsection{台长}
\begin{itemize}
\item Ira Sprague Bowen，1948–1964
\item Horace Welcome Babcock，1964–1978
\item Maarten Schmidt，1978–1980
\item Gerry Neugebauer，1980–1994
\item James Westphal，1994–1997
\item Wallace Leslie William Sargent，1997–2000
\item Richard Ellis，2000–2006
\item Shrinivas Kulkarni，2006–2018
\item Jonas Zmuidzinas，2018–2025
\item Mansi Kasliwal，2025–$^{[14]}$
\end{itemize}
\subsubsection{帕洛玛天文台与光污染}
南加州周边区域的大部分地区已采用遮光照明方式，以减少可能影响天文台的光污染$^{[15]}$。
\subsection{望远镜与仪器}
\begin{figure}[ht]
\centering
\includegraphics[width=6cm]{./figures/0287c69c9b783306.png}
\caption{黑尔望远镜穹顶} \label{fig_Haumea_3}
\end{figure}
\begin{figure}[ht]
\centering
\includegraphics[width=6cm]{./figures/5ac02f531ad92d9d.png}
\caption{黑尔望远镜的组成部件} \label{fig_Haumea_4}
\end{figure}